\answer{OA31}
用长度为 $N$ 的循环缓冲区保存窗口，以便在满窗时取出最旧值。设当前窗口大小为 $k=\min(t,N)$，维护
\[
 S=\sum_{\text{window}}x_i,\qquad Q=\sum_{\text{window}}x_i^2.
\]
未满窗时，新值 $b$ 到来后令 $S'=S+b$、$Q'=Q+b^2$、$k'=k+1$。满窗时设移出的值为 $a$，则
\[
 S'=S-a+b,\qquad Q'=Q-a^2+b^2,
 \qquad k'=N.
\]
输出
\[
 \mu'=\frac{S'}{k'},\qquad
 v'=\frac{Q'}{k'}-(\mu')^2.
\]
恒等式 $\sum(x_i-\mu)^2=Q-k\mu^2$ 证明了输出正确。每次只需常数次算术操作，故在单位代价算术模型下单次更新时间为 $O(1)$；缓冲区占 $O(N)$ 空间。$N=1$ 时方差恒为零。

当数值很大而波动很小时，$Q/k-\mu^2$ 会有浮点消去误差。可以改维护中心平方和 $M=\sum(x_i-\mu)^2$。满窗的更新为
\[
 \begin{aligned}
 \mu'&=\mu+\frac{b-a}{N},\\
 M'&=M+(b-a)(b+a-\mu-\mu'),\qquad v'=M'/N.
 \end{aligned}
\]
第二式由 $M=Q-N\mu^2$ 和 $N(\mu'-\mu)=b-a$ 展开即得。未满窗可用增量公式：令 $\delta=b-\mu$，则
\[
 \mu'=\mu+\frac{\delta}{k+1},\qquad
 M'=M+\delta(b-\mu').
\]
空窗的第一个值直接初始化为 $\mu=b,M=0$。中心更新减少了反复相减两个大二阶矩的风险，但浮点误差仍可能累积，必要时可从缓冲区重新计算。若需要分母为 $k-1$ 的样本方差，在 $k>1$ 时输出 $M/(k-1)$；是否无偏还取决于抽样模型。
